
\subsubsection{5.1 H-E1: Pareto Trade-offs Exist}

\textbf{Main Finding:} IF and FastIF exhibit metric crossings at all five tested compute budgets. IF achieves higher rank preservation (rho\_r), while FastIF achieves higher magnitude fidelity (rho\_m).

\begin{table*}[htbp]
\centering
\resizebox{\dimexpr 2\columnwidth+\columnsep\relax}{!}{%
\begin{tabular}{lllllll}
\toprule
Budget & IF rho\_r & FastIF rho\_r & rho\_r Diff & IF rho\_m & FastIF rho\_m & rho\_m Diff \\
\midrule
10 & 0.412 & 0.236 & +0.176 & 0.285 & 0.317 & -0.032 \\
25 & 0.487 & 0.228 & +0.259 & 0.291 & 0.322 & -0.031 \\
50 & 0.541 & 0.226 & +0.315 & 0.298 & 0.328 & -0.030 \\
75 & 0.583 & 0.224 & +0.359 & 0.302 & 0.331 & -0.029 \\
100 & 0.618 & 0.226 & +0.392 & 0.306 & 0.333 & -0.027 \\
\bottomrule
\end{tabular}
}
\end{table*}

The differences show opposite signs: IF outperforms FastIF on rank preservation (positive rho\_r diff) while FastIF outperforms IF on magnitude fidelity (negative rho\_m diff).

\textbf{Gate Result:} PASS (5 crossings, threshold was at least 2)

Additional observations from the H-E1 experiments:
\begin{itemize}
\item TRAK showed increasing correlation with compute budget (rho\_r from 0.40 at budget 10 to 0.57 at budget 100)
\item TracIn achieved high correlation values (rho\_r = 1.000) because it uses the same gradient computation approach as the ground truth proxy
\end{itemize}

\subsubsection{5.2 H-M1: Convex Settings Show Metric Coupling}

\textbf{Main Finding:} In logistic regression, cross-metric partial correlations exceed 0.99 at all compute levels.

\begin{table}[htbp]
\centering
\resizebox{\columnwidth}{!}{%
\begin{tabular}{lll}
\toprule
Budget & corr(rho\_r, rho\_m given budget) & Status \\
\midrule
10 & 0.9961 & PASS \\
25 & 0.9945 & PASS \\
50 & 0.9899 & PASS \\
75 & 0.9905 & PASS \\
100 & 0.9916 & PASS \\
\bottomrule
\end{tabular}
}
\end{table}

Convexity was verified with positive-definite Hessian (eigenvalues ranging from 0.01 to 0.03). IF with exact Hessian inverse achieved near-perfect correlation (rho\_r = rho\_m = 0.9999).

\textbf{Gate Result:} PASS (minimum correlation 0.9899 >= 0.95 threshold)

\subsubsection{5.3 H-M2: Deep Networks Show Metric Decoupling}

\textbf{Main Finding:} In ResNet-18, R-squared from regressing metrics on approximation error norm drops substantially compared to convex settings.

\begin{table}[htbp]
\centering
\resizebox{\columnwidth}{!}{%
\begin{tabular}{llll}
\toprule
Setting & R-squared (rho\_r) & R-squared (rho\_m) & R-squared (avg) \\
\midrule
Convex & 0.269 & 0.160 & 0.214 \\
Deep & 0.062 & 0.007 & 0.034 \\
Delta & -77\% & -96\% & -84\% \\
\bottomrule
\end{tabular}
}
\end{table}

Cross-metric correlations in deep networks varied widely across budgets, ranging from -0.45 to 0.99, compared to consistently high correlations in convex settings.

\textbf{Gate Result:} PASS (R-squared = 0.034 < 0.80 threshold)

\textbf{Note:} The convex R-squared (0.214) was lower than might be expected from theory. This occurred because IF achieved near-perfect correlation (error norm approximately 6.3) while gradient-based methods (TRAK, TracIn, FastIF) had higher error norms (approximately 680-693), creating a bimodal distribution. The relative drop from convex to deep settings remains substantial.

\subsubsection{5.4 H-M3: Methods Identify Different Influential Examples}

\textbf{Main Finding:} Different methods share less than 1\% of their top-50 influential examples.

\begin{table}[htbp]
\centering
\resizebox{\columnwidth}{!}{%
\begin{tabular}{llll}
\toprule
Budget & Min Jaccard & Mean Jaccard & Disagreement \\
\midrule
10 & 0.0034 & 0.0056 & 99.7\% \\
25 & 0.0032 & 0.0047 & 99.7\% \\
50 & 0.0026 & 0.0042 & 99.7\% \\
75 & 0.0041 & 0.0061 & 99.6\% \\
100 & 0.0024 & 0.0051 & 99.8\% \\
\bottomrule
\end{tabular}
}
\end{table}

The minimum Jaccard similarity of 0.0024 was observed between TracIn and FastIF, both of which are gradient-based methods. This suggests that even methods from the same paradigm family can identify different influential examples due to implementation differences.

Cross-paradigm and same-paradigm method pairs showed approximately equal mean Jaccard (0.0052 vs 0.0051), indicating that paradigm classification does not predict method agreement.

\textbf{Gate Result:} PASS (Jaccard = 0.0024 < 0.70 threshold)

\subsubsection{5.5 Summary}

\begin{table*}[htbp]
\centering
\resizebox{\dimexpr 2\columnwidth+\columnsep\relax}{!}{%
\begin{tabular}{lllll}
\toprule
Hypothesis & Gate Type & Predicted & Observed & Status \\
\midrule
H-E1 (Existence) & MUST\_WORK & >= 2 crossings & 5 crossings & PASS \\
H-M1 (Convex Coupling) & MUST\_WORK & corr >= 0.95 & 0.9899 & PASS \\
H-M2 (Deep Decoupling) & MUST\_WORK & R-squared < 0.80 & 0.034 & PASS \\
H-M3 (Disagreement) & SHOULD\_WORK & Jaccard < 0.70 & 0.0024 & PASS \\
\bottomrule
\end{tabular}
}
\end{table*}
